% ==========================
% Chapter 5 - reading the surface across age
% rib chapter (about ~12pp). It follows up on what was promised in Ch3: age taken back up. 
% claim: separability, aka age its own axis because vuln score is age-invariant by construction, so the split recovers a gradient the profile can't see
 
% §1 "if we want X, problems": want who, map standardizwe age away
% 2: Why age is its own axis - separability. DA-PCA invariant across bands
% §3 THe four age groups - 4 bands, 165 units./ Point back to CH3 identification 
% §4 the gradient - sim-bracketed illustration; flagged not real yet
% §5 what the reading leaves open - five beat close, hand to validation by what it does 
 
 
\chapter{Reading the surface across age}\label{chapter:agefork}
 
\section{Introduction}
 
The map read heat--mortality risk with age standardized away. Standardization was the right choice for reading the neighbourhood: it set every neighbourhood to a common age structure, so that what varied across the map was the place not its age. But the same step that isolates the neighbourhood discards the age structure it standardized over, and age is not a nuisance here.
 
This chapter takes age back up. We split the design by age group and read the heat--mortality association separately within each, recovering the gradient the standardized map combined into one number.
 
\section{Why age is its own axis}
 
A neighbourhood's vulnerability profile does not change with age. The profile is built from what a place measurably is---its population structure, its housing, its greenness---and those inputs are the same whether we read the risk for the youngest band or the oldest. The principal components that summarize the profile are therefore identical across the four age groups: one score per neighbourhood, carried unchanged into every band. 
 
This is what makes age a separate axis. If age modified risk only
through the profile, the split would recover nothing the map had not
already shown: the age groups would differ only where the profile
differed, and standardizing age away would have lost no signal worth
recovering.
 
 
\section{The four age groups}
 
We split the population into four age groups: under 65, 65 to 74, 75 to
84, and 85 and over. The groups are cut where the mortality rate turns---
nearly flat below 65, then rising steeply enough that each decade past it
carries a different risk of dying in the heat. The oldest group is set
apart because its rate runs several times that of the band below.
% [PH: "several times" is the simulated slope; replace with real 85+/75-84 rate ratio when CVSD lands.]
 
Four groups across 41 metropolitan areas give 164 units to pool, and that
164 units identify a five-parameter curve was settled in
Chapter~\ref{chapter:nationalmap}. What the split changes is the unit.
Each unit is now a city read at one age, not a city, and the pooled
association is free to bend across the age groups rather than hold to one
national shape.
 
\section{The gradient}
 
With the unit redefined, the second stage estimates a separate pooled
curve for each of the four age groups, and downscaling reads each curve at
every dissemination area in turn. The output is four national surfaces
where the map held one: the same neighbourhood now carries a
heat--mortality risk for the under-65 population and a different one for
the 85-and-over, and the difference between them is the age gradient the
standardized map combined away.
 
On simulated mortality, the pipeline recovers a known gradient. The
simulation was built with an age slope set in advance: a baseline
mortality rate that climbs across the four bands. The fitted surfaces
return that climb. Simulated deaths rise from the youngest group to the
oldest by roughly an order of magnitude, in the proportion the simulation
specified.
% [SIM ONLY: Toronto sim deaths 10,233 (0–64) → 77,210 (85+), 5 warm-seasons,
%  seed 42. Illustration of recovery, NOT a finding.]
This is a recovery check, not a result. It shows the split reads an age
gradient when one is present, and says nothing about the gradient Canada
actually carries.
 
The real gradient enters with the mortality data. The curves in this
chapter are fitted to simulated counts; the magnitude of the age
difference, and whether the oldest group's curve bends more steeply into
the heat than the youngest, are quantities the real run returns and this
draft leaves open.
% [Blocker: real per-band gradient — magnitude + old-vs-young steepness —
%  pending CVSD. Populate the four-surface comparison here.]
 
\section{What the reading leaves open}
 
There is, of course, no age gradient read here that the data identified on
its own. Each band's curve is read off the pooled national model, the same
borrowing the map rested on, and the oldest group leans on it hardest.
85-and-over is the thinnest band in every city, so the age read is most
borrowed exactly where the risk is highest. The split does not lift this
limit; it carries it, one curve per age group.
 
What the split does deliver is the age axis itself. The standardized map
held a neighbourhood's risk to a single number; reading it across the four
groups recovers the difference that number combined away, and lets the
risk for the oldest population be seen apart from the youngest rather than
averaged into it. The map left one question open: whether a risk read off
the national pooling holds at the scale of the single neighbourhood. The
age gradient inherits that question. The next chapter puts the pooled
surface to the test against risk estimated locally, and the age split is
tested on the same terms.